\documentclass{article}

\PassOptionsToPackage{numbers, compress}{natbib}
\usepackage[main, final]{neurips_2026}

\usepackage[utf8]{inputenc}
\usepackage[T1]{fontenc}
\usepackage{hyperref}
\usepackage{url}
\usepackage{booktabs}
\usepackage{amsmath,amssymb}
\usepackage{graphicx}
\usepackage{float}
\usepackage{enumitem}
\usepackage{multirow}
\usepackage{makecell}
\usepackage{tabularx}

\title{Designing a Domain-Specific Agent Runtime:\\
Lessons from Claude Code, Codex, and Cursor}

\author{
FoldAlpha Engineering \\
\texttt{team@foldalpha.com}
}

\begin{document}

\maketitle

\begin{abstract}
We describe the redesign of the agent runtime powering FoldAlpha (\url{https://app.foldalpha.com}), a production financial analysis service, from a domain-hardcoded planner architecture to a lightweight single-loop agent, informed by architectural patterns observed in commercial coding agents (Claude Code, Codex CLI, Cursor, Windsurf).
The legacy system relied on regex-based intent classification, per-turn full-state serialization, and hardcoded loop-control heuristics --- an approach that delivered functional results but diverged fundamentally from how modern agent systems are built.
The redesigned runtime adopts a minimal single-loop architecture with three key design choices: system-prompt-level context pre-loading, a reduced tool surface, and lightweight failure recovery.
In head-to-head evaluation across 6 single-turn queries and 11 multi-turn conversation turns, the new architecture achieves 2.4$\times$ lower average latency (12.7s vs.\ 30.0s), 100\% vs.\ 91\% success rate, and eliminates catastrophic loop failures observed in the legacy system.
We argue that for domain-specific agents, \emph{context density at the system prompt level} is a more effective design lever than orchestration complexity.
\end{abstract}

%---------------------------------------------------------------------
\section{Introduction}

The emergence of LLM-powered coding agents --- Claude Code~\cite{claudecodeleak2026}, OpenAI Codex CLI~\cite{codexcli2025}, Cursor~\cite{cursorprompt2025}, Windsurf~\cite{windsurfprompt2024}, and Devin~\cite{devin2025} --- has established a new paradigm for autonomous tool-using systems.
A growing body of evidence suggests that the \emph{harness} (the orchestration layer surrounding the LLM) is where the competitive differentiation lies, not the underlying model itself~\cite{effectiveharnesses2025}.

Our team operates FoldAlpha (\url{https://app.foldalpha.com}), a production financial analysis service where users ask natural-language questions and receive data-grounded answers backed by SQL queries, news retrieval, and portfolio analysis.
The legacy production system was built as a \textbf{domain-hardcoded planner}: a multi-step loop with regex-based intent classification, explicit completion checklists, and per-turn full-state serialization.
While functional, this architecture had grown increasingly difficult to maintain and extend, and its design bore little resemblance to the patterns emerging in successful commercial agents.

This paper describes how we redesigned the runtime as a \textbf{lightweight single-loop agent}, drawing on architectural principles observed in commercial systems. We present the two architectures, compare their performance, and extract design lessons for domain-specific LLM agents.

\subsection{Contributions}

\begin{enumerate}[leftmargin=*,nosep]
  \item A structural comparison of our legacy domain-hardcoded planner against the architectural patterns of commercial coding agents, identifying fundamental design divergences.
  \item A redesigned runtime based on the single-loop agent pattern with system-prompt-level context injection, demonstrating 2.4$\times$ latency improvement and higher reliability.
  \item Design principles for domain-specific agent harnesses, grounded in both our empirical results and the observed practices of commercial systems.
\end{enumerate}

%---------------------------------------------------------------------
\section{Commercial Agent Architectures}

\subsection{Landscape}

Recent source code leaks and open-source releases have provided unprecedented visibility into how production coding agents are built. We surveyed five major systems, summarized in Table~\ref{tab:agents}.

\begin{table}[t]
\centering
\caption{Architectural comparison of commercial coding agents.}
\label{tab:agents}
\small
\begin{tabular}{@{}lllll@{}}
\toprule
\textbf{Dimension} & \textbf{Claude Code} & \textbf{Codex CLI} & \textbf{Cursor} & \textbf{Windsurf} \\
\midrule
Loop type & Single + delegation & Stateless single & Dual (Plan/Exec) & Single \\
Tool count & 40+ & 100+ skills & 15+ & 9 \\
Implementation & TypeScript/Bun & Rust & IDE plugin & IDE plugin \\
Context strategy & 5 compaction strategies & Prefix caching & Symbol tracing & Persistent memory \\
Intent classification & None & None & None & None \\
Pre-planning step & None & None & Separate mode & None \\
Failure recovery & Retry + compaction & Stateless retry & Edit confidence & Proactive execution \\
\bottomrule
\end{tabular}
\end{table}

\subsection{Key Observations}

Three patterns are consistent across all surveyed systems:

\textbf{(1) Simple loop, no intent classification.}
Every agent uses a single-loop architecture where the LLM autonomously decides the next action at each step. None employ regex-based or classifier-based intent routing. The LLM itself serves as the router.

\textbf{(2) Context management over orchestration complexity.}
The primary engineering investment is in \emph{what context reaches the LLM and how}, not in multi-stage pipelines.
Claude Code tracks 14 cache-break vectors and implements 5 compaction strategies~\cite{claudecodeleak2026}.
Codex CLI maintains stateless prefix caching where each request preserves the previous prompt as an exact prefix to maximize cache hits~\cite{codexloop2025}.

\textbf{(3) Lightweight failure recovery.}
Rather than preventing failures through complex pre-planning, these systems detect and recover from failures cheaply: retrying on empty responses, blocking repeated tool calls, and compacting context when it grows too large.

%---------------------------------------------------------------------
\section{Architecture Comparison: Legacy vs.\ Redesigned}

\subsection{Legacy: Domain-Hardcoded Planner}

Our production planner (Node.js/TypeScript, ${\sim}$700 lines) employed the following design:

\textbf{Intent classification.}
On every question, regex patterns classified the query into one of 6 categories (\texttt{compare}, \texttt{portfolio-review}, \texttt{screening}, \texttt{outlook}, \texttt{concept}, \texttt{analysis}), with 7 must-cover boolean flags (news, earnings, valuation, etc.) and a dynamically generated completion checklist.

\textbf{Per-turn state serialization.}
Every LLM call received a 2-message payload (system + user) where the user message contained the \emph{entire} accumulated state: question, conversation history, portfolio, previous result JSON, intent classification output, schema catalog (${\sim}$3K tokens), full tool history with results, current SQL, and current dataset preview. This produced a constant-but-large token payload per step.

\textbf{Hardcoded loop control.}
The system tracked repeated tool calls via loop keys (\texttt{tool\_name:args}) and outcome keys (\texttt{tool\_name:args:result\_shape}), blocking after 4 identical calls or 3 identical outcomes. A safety cap of 12 steps provided the outer bound. Early-exit heuristics triggered for concept questions after detecting sufficient evidence.

\textbf{Post-processing pipeline.}
Final responses passed through \texttt{stripMarkdownTables}, \texttt{buildStateBackedAssistantMessage} (injecting news/financial summaries from state), and an optional LLM-based display renderer for chart/table specification.

\subsection{Redesigned: Lightweight Single-Loop Agent}

The new runtime (Python/FastAPI) follows the single-loop pattern observed in commercial agents:

\textbf{No intent classification.}
The LLM receives the question, tools, and context, and autonomously decides whether to call a tool or produce a final answer. There are no regex classifiers, completion checklists, or must-cover flags.

\textbf{System-prompt-level context.}
Domain schema (${\sim}$3K tokens) is injected once into the system prompt at agent initialization. This placement enables API-level prefix caching across turns (supported by Gemini, Anthropic, and OpenAI). A lightweight state context section is appended to the system message each turn, summarizing loaded portfolio, recent tool history (last 4 calls), and last SQL result metadata.

\textbf{Conversation replay with condensation.}
The full event log (user messages, assistant responses, tool calls, observations) is replayed as the message history. An LLM-based condenser triggers after 24 events, summarizing older events to bound context growth.

\textbf{Minimal failure recovery.}
Three mechanisms provide robustness without architectural complexity:
\begin{enumerate}[nosep,leftmargin=*]
  \item \textbf{Empty response retry}: One retry when the LLM returns no content and no tool calls.
  \item \textbf{Loop detection}: Blocking identical tool+args combinations after 3 repetitions.
  \item \textbf{Response post-processing}: Stripping markdown tables when a data panel will display the SQL result.
\end{enumerate}

\textbf{Reduced tool surface.}
Three tools (\texttt{run\_sql}, \texttt{search\_news}, \texttt{get\_portfolio}) replace the legacy system's 9 tools (which included separate schema catalog, metric resolver, reporting grain resolver, price field resolver, derived metric resolver, and KR/US stock ticker resolvers). The schema guide, previously requiring a tool call to load, is now always present in the system prompt.

Table~\ref{tab:arch} summarizes the structural differences.

\begin{table}[t]
\centering
\caption{Structural comparison of legacy and redesigned architectures.}
\label{tab:arch}
\small
\begin{tabular}{@{}lll@{}}
\toprule
\textbf{Aspect} & \textbf{Legacy Planner} & \textbf{Redesigned Agent} \\
\midrule
Loop structure & Single loop + intent routing & Single loop, LLM-driven \\
Intent classification & Regex (6 types + 7 flags) & None \\
Schema delivery & Per-turn in user message & System prompt (cached) \\
State per turn & Full serialization (${\sim}$fixed tokens) & Conversation replay (growing) \\
Tool count & 9 & 3 \\
Loop control & Hardcoded (4$\times$ block + outcome dedup) & Simple counter (3$\times$ block) \\
Post-processing & State-backed message + LLM renderer & Markdown table stripping \\
Token growth pattern & Constant but large per step & Growing but cache-friendly \\
Lines of code & ${\sim}$700 (runtime only) & ${\sim}$380 (runtime only) \\
\bottomrule
\end{tabular}
\end{table}

%---------------------------------------------------------------------
\section{Evaluation}

Both runtimes were evaluated on the same infrastructure, using the same LLM (Gemini Flash) and Database containing Korean and US equity data.

\subsection{Single-Turn Benchmark}

Six representative financial queries spanning different analysis types were evaluated.

\begin{table}[t]
\centering
\caption{Single-turn evaluation. Time in seconds. \textbf{Bold} indicates best per row.}
\label{tab:singleturn}
\small
\begin{tabular}{@{}llrrrr@{}}
\toprule
\textbf{Type} & \textbf{Query} & \makecell{\textbf{Legacy}} & \makecell{\textbf{Redesigned}} & \textbf{Speedup} \\
\midrule
Earnings    & Samsung Electronics earnings     & 28.5 & \textbf{12.6} & 2.3$\times$ \\
News        & NVIDIA recent news               & 25.2 & \textbf{8.9}  & 2.8$\times$ \\
Snapshot    & NVIDIA snapshot                  & 23.8 & \textbf{13.0} & 1.8$\times$ \\
Comparison  & Hyundai vs.\ Kia                  & 30.7 & \textbf{17.1} & 1.8$\times$ \\
Screening   & KOSPI high-CAGR low-PBR          & 42.0 & \textbf{15.9} & 2.6$\times$ \\
Lookup      & KOSDAQ market cap leader         & 29.9 & \textbf{8.9}  & 3.4$\times$ \\
\midrule
\multicolumn{2}{l}{\textbf{Mean}} & 30.0 & \textbf{12.7} & \textbf{2.4$\times$} \\
\bottomrule
\end{tabular}
\end{table}

The redesigned runtime achieves 2.4$\times$ average speedup with 100\% success rate. The legacy planner's higher latency stems from two factors: (1) per-turn schema serialization in the user message forces larger token payloads without caching benefits, and (2) additional LLM calls are consumed by SQL error recovery and semantic tool resolution.

\subsection{Multi-Turn Benchmark}

Four multi-turn scenarios (11 total turns) tested follow-up continuity, cross-entity comparison, screening-to-detail drilldown, and iterative analysis.

\begin{table}[t]
\centering
\caption{Multi-turn evaluation summary.}
\label{tab:multiturn}
\small
\begin{tabular}{@{}lrrcl@{}}
\toprule
& \makecell{\textbf{Avg Latency}\\\textbf{(s)}} & \makecell{\textbf{Success}} & \makecell{\textbf{Worst Case}} \\
\midrule
Legacy planner    & 29.5 & 10/11 (91\%)  & 96.9s, 12-step loop, 21-char answer \\
Redesigned agent  & \textbf{14.4} & \textbf{11/11 (100\%)} & 28.9s, normal execution \\
\bottomrule
\end{tabular}
\end{table}

The legacy planner's worst case is notable: on an NVIDIA 5-year revenue CAGR query, it consumed 12 steps and 96.9 seconds before producing a 21-character fallback message (``Data query results organized.'').
This occurred \emph{despite} its hardcoded loop-control mechanisms, suggesting that rule-based loop prevention has fundamental coverage gaps.
The redesigned runtime answered the same query in 8.9 seconds with a single SQL call.

Follow-up turns in both systems were generally faster than initial turns. The redesigned runtime's state snapshot mechanism correctly propagated context across turns, supporting references like ``show me the top stock from the previous screening'' without re-querying.

\subsection{General-Purpose Agent Baseline}

For additional context, we tested the same queries using Claude Code (a general-purpose coding agent) as a baseline, in two configurations:

\begin{table}[H]
\centering
\caption{Baseline comparison: domain-specific vs.\ general-purpose agents.}
\label{tab:baseline}
\small
\begin{tabular}{@{}lrrl@{}}
\toprule
\textbf{Configuration} & \textbf{Time (s)} & \textbf{Retries} & \textbf{Notes} \\
\midrule
General-purpose (no context)   & ${\sim}$180 & 4 & Environment discovery overhead \\
General-purpose (with guide)   & 73          & 0 & Pre-loaded DB access guide \\
Redesigned domain agent        & \textbf{12.6} & 0 & Schema in system prompt \\
Legacy domain planner          & 28.5        & 1 & Schema in user message \\
\bottomrule
\end{tabular}
\end{table}

The 14.4$\times$ gap between the general-purpose agent (no context) and the domain agent confirms that \textbf{pre-loaded domain context} is the primary efficiency lever, independent of agent sophistication. Providing the general-purpose agent with a database access guide reduced its latency by 2.5$\times$ and eliminated all retries, further supporting this conclusion.

\subsection{Answer Quality}

Both domain runtimes produced structured answers with markdown formatting, data-grounded analysis, and year-over-year comparisons.
The redesigned runtime generated moderately longer answers on average (1,050 vs.\ 780 characters) with more explicit interpretation.
In the screening scenario, the redesigned runtime independently noted base-effect risks in growth rate calculations --- a qualitative indicator of effective LLM reasoning within the simplified harness.

Both systems occasionally produced different but valid SQL interpretations for the same query (e.g., different minimum-value filters for growth rate screening). We consider this acceptable variance in LLM-generated SQL rather than an architectural issue.

%---------------------------------------------------------------------
\section{Discussion}

\subsection{Why Simpler Won}

The legacy planner's complexity (intent classification, completion checklists, must-cover flags, outcome deduplication) represented a significant engineering investment.
Yet this complexity provided \textbf{no measurable advantage} in our evaluation: it was slower on every query type and failed catastrophically in one case where the redesigned runtime succeeded trivially.

We attribute this to a fundamental tension: \textbf{domain heuristics compete with the LLM's own reasoning}.
The regex-based intent classifier and must-cover flags attempted to guide the LLM's behavior through structured constraints, but the LLM already possesses the capability to make these judgments when given adequate context.
By removing these constraints and instead providing richer context (the schema in the system prompt), we allowed the LLM to reason more naturally, producing both faster and more reliable results.

This parallels the commercial agent landscape: Claude Code's 512K-line harness invests in context management and failure recovery, not intent classification~\cite{claudecodeleak2026}. Codex CLI's Rust rewrite optimizes for stateless prefix caching, not planning pipelines~\cite{codexloop2025,pragmaticcodex2025}.

\subsection{System Prompt as Cached Context Surface}

A key architectural insight is that domain knowledge in the system prompt occupies a \textbf{cache-friendly position} in the API call.
Modern LLM APIs (Gemini, Anthropic, OpenAI) support prefix caching: if the beginning of the prompt is identical across requests, the provider can skip re-processing those tokens.

The legacy planner's design of embedding schema in the user message \textbf{defeats this caching} because the user message (containing accumulated tool history and dataset previews) changes every turn.
The redesigned runtime's system-prompt placement preserves the cacheable prefix, reducing both latency and API cost as conversations grow.

\subsection{Minimal Viable Defenses}

Our three failure-recovery mechanisms (empty-response retry, loop detection, response post-processing) form what we term a \textbf{minimal viable defense set}.
These same patterns appear in more sophisticated forms across all surveyed commercial agents.
We argue that for domain-specific agents, this minimal set provides sufficient robustness without the maintenance burden of elaborate heuristic systems.

\subsection{Limitations and Future Work}

Our evaluation uses a single LLM (Gemini Flash); cross-provider evaluation is needed.
The query set (17 total turns) is small; broader benchmarking would strengthen the findings.
The redesigned runtime does not yet implement intelligent display rendering (chart vs.\ table selection), which the legacy system handled via an additional LLM call.
The schema pre-loading approach adds ${\sim}$3K tokens to every system prompt, which may impact models with smaller context windows.

Future directions include: multi-provider evaluation, context condensation optimization for extended conversations, display rendering without additional LLM calls, and adaptive context injection that loads domain knowledge relevant to the detected query type.

%---------------------------------------------------------------------
\section{Conclusion}

We redesigned a production financial analysis runtime from a domain-hardcoded planner to a lightweight single-loop agent, guided by architectural patterns observed in commercial coding agents.
The redesigned system achieves 2.4$\times$ lower latency, higher success rates, and eliminates catastrophic failure modes, while reducing the codebase by 45\%.

The central finding is that \textbf{context density and placement} --- not orchestration complexity --- is the primary lever for domain-specific agent performance.
Pre-loading domain knowledge in the system prompt, reducing the tool surface to essential capabilities, and adding minimal failure recovery produces better outcomes than elaborate intent classification, completion tracking, and heuristic loop control.

This conclusion aligns with the design philosophy of leading commercial agents, which uniformly invest in context management over planning complexity.
For practitioners building domain-specific LLM agents, we recommend starting with the simplest possible loop and investing engineering effort in what context the LLM sees, not how it is told to reason.

The runtime source code, evaluation framework, and benchmark results are publicly available at \url{https://github.com/sbyoun/finllm-agent}.

%---------------------------------------------------------------------
\begin{thebibliography}{10}

\bibitem{claudecodeleak2026}
Anthropic.
Claude Code source analysis (leaked via npm source map, March 2026).
\textit{VentureBeat, The AI Corner, Alex Kim}, 2026.

\bibitem{codexcli2025}
OpenAI.
Codex CLI: Open-source coding agent.
\url{https://github.com/openai/codex}, 2025.

\bibitem{cursorprompt2025}
Cursor Agent System Prompt (Agent Prompt 2.0).
Leaked March 2025.

\bibitem{windsurfprompt2024}
Windsurf Cascade System Prompt (Tools Wave 11).
Leaked December 2024.

\bibitem{devin2025}
T.~Endo.
Agent-native development: A deep dive into Devin 2.0's technical design.
\textit{Medium}, 2025.

\bibitem{effectiveharnesses2025}
Anthropic.
Effective harnesses for long-running agents.
\textit{Anthropic Engineering Blog}, 2025.

\bibitem{codexloop2025}
OpenAI.
Unrolling the Codex agent loop.
\textit{OpenAI Blog}, 2025.

\bibitem{pragmaticcodex2025}
G.~Orosz.
How Codex is built.
\textit{The Pragmatic Engineer Newsletter}, 2025.

\end{thebibliography}

\end{document}
